\documentclass[12pt]{article}
\usepackage[T1]{fontenc}
\usepackage[utf8]{inputenc}
\usepackage[brazil]{babel}
\usepackage[margin=1in]{geometry}
\setlength{\parindent}{0pt}
\begin{document}
\section*{Tema 64}
Processo: PEDILEF 2009.72.54.005939-9/ SC\\
Situacao: Julgado\\
Ramo: DIREITO TRIBUTÁRIO\\
Questao: Saber se incide contribuição previdenciária sobre auxílio-alimentação para cargos em comissão.\\
Tese: Os empregados submetidos ao Regime Geral da Previdência Social (RGPS), aí incluídos os exercentes de cargo em comissão, de livre nomeação e exoneração, em caráter exclusivo, se sujeitam à incidência de contribuição previdenciária sobre as parcelas percebidas a título de Auxílio-Alimentação dada a sua natureza salarial, com base nos termos do art. 40, § 13º, da CF/88 c.c. art. 28, I, da Lei n. 8.212/91, salvo, se tal pagamento for “in natura”, isto é, quando a própria empresa fornece a alimentação. Vide Súmula 67 da TNU.\\
Fonte: https://www.cjf.jus.br/phpdoc/virtus/
\end{document}
